% !PRINTED 107
\weberpage{107}
On the other hand, from the formulae (8) of § 21 there follows, if we there set $v = -\dfrac{\omega}{4}$ and $= -\dfrac{1}{4}$,
% [sic] The German source omits the variable before the second equals sign; the context suggests v again.
\[
\begin{aligned}
\vartheta_{11}\left(\frac{\omega}{4}\right) &= i\vartheta_{01}\left(\frac{\omega}{4}\right) \\
\vartheta_{11}\left(\frac{1}{4}\right) &= \vartheta_{10}\left(\frac{1}{4}\right)
\end{aligned}
\]
and accordingly the two functions (12), which are linearly expressible by two $\vartheta$-squares, agree, apart from constant factors, with
\[
\vartheta_{01}^2(u) + \vartheta_{11}^2(u), \quad \vartheta_{10}^2(u) - \vartheta_{11}^2(u).
\]
The constant factors are obtained immediately, for $u = 0$, from the relations (6), (7):
\weq{13}{\vartheta_{00}\left(0, \frac{\omega}{2}\right)\vartheta_{01}\left(u, \frac{\omega}{2}\right) = \vartheta_{01}^2(u) + \vartheta_{11}^2(u),}
\weq{14}{2\vartheta_{00}(0, 2\omega)\vartheta_{10}(2u, 2\omega) = \vartheta_{10}^2(u) - \vartheta_{11}^2(u).}
From these one obtains the last two formulae by changing $u$ into $u + \dfrac{1}{2}$ and $u + \dfrac{\omega}{2}$ [or else in the same way as (13), (14)]:
\weq{15}{\vartheta_{00}\left(0, \frac{\omega}{2}\right)\vartheta_{00}\left(u, \frac{\omega}{2}\right) = \vartheta_{00}^2(u) + \vartheta_{10}^2(u),}
\weq{16}{2\vartheta_{00}(0, 2\omega)\vartheta_{00}(2u, 2\omega) = \vartheta_{00}^2(u) + \vartheta_{01}^2(u).}
From these, manifold relations between the zero values of the $\vartheta$-functions can be derived, of which we adduce only the three following ones. The first two of them flow from (8), (9), while the last follows from the first equation (11) if one replaces $\omega$ by $\omega{:}2$, assumes $u = 1/4$, and takes into account that $\vartheta_{11}(1/4) = \vartheta_{10}(1/4)$.
\weq{17}{\begin{aligned}
\sqrt{\vartheta_{00}\vartheta_{01}} &= \vartheta_{01}(0, 2\omega), \\
\sqrt{\vartheta_{00}\vartheta_{10}} &= \frac{1}{\sqrt{2}}\vartheta_{10}\left(0, \frac{\omega}{2}\right), \\
\sqrt{\vartheta_{01}\vartheta_{10}} &= \vartheta_{10}\left(\frac{1}{4}, \frac{\omega}{2}\right).
\end{aligned}}
These formulae are of interest because they represent the square roots as single-valued functions of $\omega$.
